% Body-only note. The parent report supplies amsmath, amssymb, amsthm,
% booktabs and hyperref. All results are ordinary mathematics, not kernel proofs.
\section{A smaller finite formalism by interpreting raw assignments}
\label{sec:selective-equality}

The previous completed FC15$^*$ formalism retained congruence for the two
assignment functions $V,C$, for membership $E$, and for incidence
$\mathsf{Inc}$. Its strongest delivered reduction had 20 universal clauses,
53 literals and 1,134 flat template bits. That argument reconstructed a
native assignment/class model on an equality quotient.
Here we prove a stronger reduction by constructing only the set-theoretic
reduct needed for the consistency comparison. Raw class codes and raw
assignments remain outside the quotient interpretation. This removes
all update congruence, all incidence congruence, and lookup congruence in
the selector argument.

The main named-selector target has \textbf{19 clauses, 45 literals, and
971 template bits}. Its extensional control has \textbf{20 clauses,
49 literals, and 1,138 bits}. Both use two explicitly encoded constants,
nine proper functions and four ordinary binary predicate symbols,
including the equality surrogate. These are exact formalism sizes,
not Turing-machine state counts. The second constant is charged in the
template and ground-term encodings. There is no hidden initial tape.
Sections~\ref{se:fusion}--\ref{se:witness-fusion} then reduce the proper
function count to eight at no payload cost, to six for 40 extra bits,
or to five for 56 extra bits, by explicit two-coordinate model constructions.

\subsection{The exact language and four additional clauses}

The constants are $d,q$. The proper functions are
\[
 V:2,\quad C:3,\quad N:2,\quad A:2,\quad Q:2,
 \quad W:3,\quad Y:3,\quad Z:3,\quad H:2.
\]
The binary predicates are $E,\mathsf{Inc},G,\sim$. In the finite theory
$\sim$ is an ordinary predicate, not built-in logical equality.
Every variable below is universally quantified. Write
\[
 R(b,t,x,y):=\mathsf{Inc}(C(C(b,d,x),q,y),t),\qquad
 D(x):=W(C(d,q,x),d,A(d,q)).
\]
The fifteen inherited native clauses, with their pseudo-equalities now
read as $\sim$, are exactly
\begin{align*}
 &V(C(s,j,z),j)\sim z,\\
 &(k\sim j)\lor (V(C(s,j,z),k)\sim V(s,k)),\\
 &\mathsf{Inc}(s,N(a,b))\lor\mathsf{Inc}(s,a),\\
 &\mathsf{Inc}(s,N(a,b))\lor\mathsf{Inc}(s,b),\\
 &\neg\mathsf{Inc}(s,N(a,b))\lor\neg\mathsf{Inc}(s,a)
        \lor\neg\mathsf{Inc}(s,b),\\
 &\neg\mathsf{Inc}(s,A(j,k))\lor E(V(s,j),V(s,k)),\\
 &\mathsf{Inc}(s,A(j,k))\lor\neg E(V(s,j),V(s,k)),\\
 &\neg\mathsf{Inc}(C(s,j,z),a)\lor\mathsf{Inc}(s,Q(j,a)),\\
 &\neg\mathsf{Inc}(s,Q(j,a))
        \lor\mathsf{Inc}(C(s,j,W(s,j,a)),a),\\
 &G(b,t)\lor R(b,t,H(b,t),Y(b,t,u)),\\
 &G(b,t)\lor E(Z(b,t,u),Y(b,t,u)),\\
 &G(b,t)\lor\neg E(Z(b,t,u),u),\\
 &\neg G(b,t)\lor\neg E(x,H(b,t))\lor\neg R(b,t,x,y)
        \lor E(y,H(b,t)),\\
 &\neg G(b,t)\lor E(D(x),x)\lor\neg R(b,t,x,y)
        \lor E(y,H(b,t)),\\
 &\neg G(b,t)\lor\neg E(y,H(b,t))\lor R(b,t,Y(b,t,y),y).
\end{align*}
All abbreviations are expanded in the emitted JSON and bit streams.
The source44 upper construction chooses the old selector as $N(d,d)$;
our new language leaves $q$ independent and requires only the separation
needed by the lower proof.

Add precisely these four clauses:
\begin{align}
 &\neg(x\sim y)\lor\neg(x\sim z)\lor (y\sim z),
          \tag{SE1}\label{se:eq}\\
 &\neg(s\sim t)\lor (V(s,j)\sim V(t,j)),
          \tag{SE2}\label{se:v}\\
 &\neg(a\sim b)\lor\neg(x\sim y)\lor\neg E(a,x)\lor E(b,y),
          \tag{SE3}\label{se:mem}\\
 &\neg(d\sim q).
          \tag{SE4}\label{se:sep}
\end{align}
Call the resulting theory $\mathsf{SE19}$. Clause \eqref{se:v} transports
only the first argument of $V$, with the literally same selector $j$.
There is no transport clause for $C,N,A,Q,W,Y,Z,H$, for $G$ or
$\mathsf{Inc}$, or for the second argument of $V$.

\subsection{An equality quotient with sufficiently many selector classes}

Let $M$ be any set-sized model of these nineteen clauses, with raw carrier
$U$. The first native clause at $s=j=d,z=x$ supplies
$p_x\sim x$, where $p_x=V(C(d,d,x),d)$. Applying \eqref{se:eq} to
$(p_x,x,x)$ proves $x\sim x$. Reflexivity and \eqref{se:eq} give symmetry;
symmetry and \eqref{se:eq} give transitivity. Thus $\sim$ is an
equivalence relation even though no reflexivity clause was emitted.

Put $D=U/{\sim}$. We use brackets for quotient classes and interpret
\[
 E_D([x],[y])\quad\Longleftrightarrow\quad E_M(x,y).
\]
Clause \eqref{se:mem}, together with symmetry, proves this is well-defined
in both arguments. No function other than this quotient relation will
be interpreted on $D$.

For raw objects put
\[
 P(x,y)=C(C(d,d,x),q,y).
\]
The first two native clauses and \eqref{se:sep} imply
\[
 V(P(x,y),d)\sim x,\qquad V(P(x,y),q)\sim y.
\]
If $P(x,y)\sim P(x',y')$, clause \eqref{se:v} at $d$ and $q$
therefore implies $x\sim x'$ and $y\sim y'$. This is a cancellation
statement; it does not assume that $P$ itself descends to a quotient
function. Define raw terms
\[
 a_0=P(q,d),\qquad a_{n+1}=P(d,a_n).
\]
They belong to pairwise distinct equivalence classes. An equality involving
$a_0$ and a later term contradicts the first projection and $d\not\sim q$;
an equality between two later terms reduces by the second projection to
one with smaller indices. Thus every required finite number of selector
classes avoiding $[d],[q]$ exists. Equivalently, after choosing raw
representatives, the cancellation property gives an injection $D^2\to D$.
The finite argument already suffices for formula compilation.

\subsection{The truth lemma that removes the remaining congruence}

Choose raw selector representatives $j_v$ for the finitely many variable
names of a fixed equality-free $E$-formula, after renaming bound variables,
such that distinct names have distinct $\sim$-classes. For each raw
assignment $s\in U$, define the finite environment
\[
 \eta_s(v)=[V(s,j_v)]\in D.
\]
The assignment clauses say exactly that replacing $s$ by
$C(s,j_v,z)$ updates this finite environment at $v$ to $[z]$ and
preserves every other selected coordinate. This assertion uses the
actual raw update object. It does not identify updates of equivalent
arguments, and it does not transport incidence along equality.

Build a raw class object $t_\varphi$ by recursion:
\begin{align*}
 t_{E(v,w)}&=A(j_v,j_w),\\
 t_{\operatorname{NAND}(\varphi,\psi)}&=N(t_\varphi,t_\psi),\\
 t_{\exists v\,\varphi}&=Q(j_v,t_\varphi).
\end{align*}
Other classical connectives and falsity are represented by their ordinary
finite NAND formulas. The native clauses prove the exact raw contracts
\begin{align*}
 \mathsf{Inc}(s,N(a,b))&\leftrightarrow
   \neg(\mathsf{Inc}(s,a)\land\mathsf{Inc}(s,b)),\\
 \mathsf{Inc}(s,A(j,k))&\leftrightarrow E(V(s,j),V(s,k)),\\
 \mathsf{Inc}(s,Q(j,a))&\leftrightarrow
   \exists z\in U\;\mathsf{Inc}(C(s,j,z),a).
\end{align*}
The last equivalence uses the supplied $W$ witness in its nontrivial
direction. Consequently structural induction gives, for every raw $s$,
\begin{equation}
 \mathsf{Inc}_M(s,t_\varphi)
 \quad\Longleftrightarrow\quad
 (D,E_D)\models\varphi[\eta_s].
 \label{se:truth}
\end{equation}
For the quantifier step, the induction hypothesis applies to each actual
raw object $C(s,j_v,z)$; its environment is the required update, and
the map $z\mapsto[z]$ is onto $D$. This deals with all raw witnesses,
not merely a selected transversal. It is the point at which attempting
to reconstruct a quotient update function would impose unnecessary
congruence conditions.

No arbitrary class is claimed to be invariant under equivalent raw
assignments or equivalent raw class codes. Equation \eqref{se:truth}
establishes the required invariance only for the class objects constructed
from finite formulas, and that is enough.

\subsection{All parameters and the exact supported-closure interface}

Fix an arbitrary equality-free matrix $\varphi(x,y,\bar p)$ in $E$.
Assign $d,q$ to the free input/output names, with all parameter and bound
names in other pairwise distinct selector classes. Finite raw $C$
updates produce an assignment $b$ realizing any prescribed tuple of
parameter classes in $D$. The truth lemma gives a raw class $t$ with
\begin{equation}
 R_M(b,t,u,v)
 \quad\Longleftrightarrow\quad
 (D,E_D)\models\varphi([u],[v],\bar p)
 \quad(u,v\in U).
 \label{se:r}
\end{equation}
Thus the selected relation $R_M$ factors through both object arguments,
although neither $C$ nor $\mathsf{Inc}$ is required to do so separately.

Write $\operatorname{Empty}(x)$ for $\forall z\,\neg E_D(z,x)$
and $y\subseteq_E u$ for $\forall z(E_D(z,y)\to E_D(z,u))$.
The raw EPSC instance required from the inherited interface is
\begin{align*}
 &[\forall x\,\exists u\,\forall y
       (\varphi(x,y,\bar p)\to y\subseteq_E u)]\\
 &\qquad\Longrightarrow\quad
 \exists c\bigl[
 \forall x,y\bigl(((\operatorname{Empty}(x)\lor E_D(x,c))
     \land\varphi(x,y,\bar p))\to E_D(y,c)\bigr)\\
 &\hspace{47mm}\land\forall y(E_D(y,c)\to\exists x\,
                  \varphi(x,y,\bar p))\bigr].
\end{align*}

Suppose its antecedent holds in $D$. If the raw flag $G_M(b,t)$ were
false, apply local subset boundedness at $[H(b,t)]$ and choose a raw
representative $u$ of the supplied bound. Clauses 10--12 give an output
$Y(b,t,u)$ satisfying $R$, a member $Z(b,t,u)$ of that output, and
nonmembership of the same object in $u$. Equation \eqref{se:r} and
the definition of $E_D$ contradict the chosen bound. Hence $G_M(b,t)$
is true.

Take $c=[H(b,t)]$. Clause 13 gives recurrent closure. If $[x]$ is
empty in $D$, then in particular $E_M(D(x),x)$ is false, whatever value
the raw term $D(x)$ has; clause 14 gives the empty-input seed. Finally
clause 15 supplies the supporting input $[Y(b,t,y)]$ for each hull
member $[y]$. The exact support and all-empty-object seed are retained.
This proves every raw EPSC instance, with unrestricted finite parameters
and quantifier complexity, in the quotient structure $(D,E_D)$.

\subsection{The extensional control}

Let $\operatorname{Xor}(a,b)$ denote the inherited NAND class term
\[
 N(N(a,N(a,b)),N(b,N(a,b))).
\]
Add the single clause
\[
 (j\sim d)\lor(j\sim q)\lor(a\sim b)\lor
 \mathsf{Inc}\bigl(P(a,b),
 Q(j,\operatorname{Xor}(A(j,d),A(j,q)))\bigr).
\]
The resulting theory is $\mathsf{SE20}$. For any two distinct classes
$[a],[b]$, choose a selector class $[j]$ avoiding $[d],[q]$.
The projection identities and the raw truth lemma identify the final
literal with
\[
 \exists z\in D\;[E_D(z,[a])\mathbin{\mathrm{xor}}E_D(z,[b])].
\]
Therefore $E_D$ is extensional. All colliding selector cases remain
covered by the displayed guards. This derives actual Extensionality
on the quotient, without transporting the values of $N,A,Q,C$
or $\mathsf{Inc}$ through that quotient.

\subsection{Consistency calibration and its exact dependency}

Inside ZFC, the preceding constructions are set-sized quotient and
satisfaction arguments. First-order completeness applied to
$\operatorname{Con}(\mathsf{SE19})$ produces an initial set model;
its quotient gives a model of full raw equality-free EPSC. The inherited
ProveIt raw-EPSC calibration, which uses a hereditary-saturation
interpretation, gives
\[
 \operatorname{Con}(\mathsf{SE19})\to
 \operatorname{Con}(\mathrm{ZFC}).
\]
For $\mathsf{SE20}$ the quotient is extensional, so the inherited
direct Extensionality-plus-EPSC bootstrap gives ZF without Foundation
or Choice, followed by the existing Foundation/constructibility
consistency transports. This is an alternative calibration dependency
that does not require the raw hereditary-saturation interpretation.

Conversely, the existing countable ZFC model construction of native
FC15$^*$/FC16$^*$ provides a model of each new finite theory. Interpret
$\sim$ as genuine identity and the new constant $q$ as the old term
$N(d,d)$. In that native model NAND proves $d\ne N(d,d)$.
All four added clauses hold under identity, and all native rows are
unchanged by this assignment to $q$. Thus, with precisely the named
inherited calibrations, ordinary ZFC proves
\[
 \operatorname{Con}(\mathsf{SE19})\leftrightarrow
 \operatorname{Con}(\mathrm{ZFC}),\qquad
 \operatorname{Con}(\mathsf{SE20})\leftrightarrow
 \operatorname{Con}(\mathrm{ZFC}).
\]
These are new mathematical reductions of the completed finite payload.
They are not claimed as new proofs of the inherited EPSC calibration,
as kernel-checked endpoints, or as PA-verified proof-code compilers.
In particular the argument does not license deleting equality rules in
the unrelated 245-state Hilbert/Tarski source calculus.

\subsection{Exact costs and an executable semidecision specification}

The supplied generator imports byte-exact source44 clause builders,
replaces each old selector term by the charged constant when selected,
and appends exactly \eqref{se:eq}--\eqref{se:sep}. Its flat stream has
three bits per clause length, three bits per literal's sign/predicate,
four bits per term node, and a three-bit terminating zero length.
Four local variable opcodes, two constants and nine function opcodes
fit in four bits. The equality predicate is included in the two-bit
predicate tag. The one-constant controls use $q=N(d,d)$ literally.

\begin{center}\small
\begin{tabular}{lrrrrr}
\toprule
Target & Clauses & Literals & Nodes & Max.\ variables & Bits\\
\midrule
Prior FC15$^*$ reduct &20&53&228&6&1134\\
$\mathsf{SE19}$, one constant &19&45&208&4&1027\\
$\mathsf{SE19}$, named $q$ &19&45&194&4&971\\
Prior FC16$^*$ reduct &21&57&276&6&1341\\
$\mathsf{SE20}$, one constant &20&49&256&4&1234\\
$\mathsf{SE20}$, named $q$ &20&49&232&4&1138\\
\bottomrule
\end{tabular}
\end{center}
The named streams save 163 and 203 bits over the prior completed
targets. Their packed byte lengths are 122 and 143, respectively,
with five and six trailing zero padding bits. Maximum term depth is
three for the raw target and five for the named extensional control.
The one-constant comparison isolates the logical reduction from the
separate constant-naming choice.

There is also an exact reduction in the grounding burden of the stated
frontend. For a frontier of $B$ terms, before deduplication, the old
20-clause target emits
\[
 B^6+7B^4+12B^3
\]
ground substitution records. The new 19-clause target emits
\[
 5B^4+13B^3+1.
\]
The extensional controls add $B^3$ to each expression. These polynomials
count all colliding substitutions too; they follow by counting local
variables in each literal template. Thus the direct exhaustive grounding
stage drops from degree six to degree four. This is a scoped compiler
benefit, not a bound on the total SAT search or the machine's halting time.

The reference refuter treats all four predicates as ordinary
propositional atom families. Its ground codec has $d=0,q=1$ and
\[
 \operatorname{code}(f_i(t_1,\ldots,t_r))
 =2+i+16\,\operatorname{tuple}(\operatorname{code}(t_1),\ldots,
                              \operatorname{code}(t_r)),
\]
with the dual Cantor pairing and seven unused unary padding
constructors. Children have smaller codes; the codec is a bijection
between naturals and the padded ground-term signature. Padding
functions are absent from the clauses and need no congruence.

At stage $n$, instantiate every local variable with every term code
$0,\ldots,n$, including collisions. Keep all resulting endpoints,
even when their codes are larger than $n$. Decide the resulting
finite ground CNF exactly; halt only on UNSAT. The stages are nested
and exhaust all ground instances. If every stage is SAT,
propositional compactness supplies a valuation of all ground atoms;
the raw term structure with that valuation is a model of the
universal clauses. Conversely every first-order model satisfies
every finite stage. Consequently the ideal refuter has the required
ZFC-provable halting-versus-inconsistency contract.

The solver adds no built-in equality or quotient-function constraints.
In particular, giving $\mathsf{Inc}$ a table on equality classes would
silently reinstate a condition that the new finite theory removed.
Resource errors propagate and never count as UNSAT. Finite-stage SAT
outputs are staging regressions, not consistency evidence. The supplied
run checks four exact payload roundtrips, twelve malformed-stream
rejections, 4,096 term-code roundtrips, 8,428 independently reconstructed
ground endpoints, and 200 small CNFs against exhaustive assignments.
At stages zero and one the raw target has 19 and 185 ground clauses;
the extensional target has 20 and 193. All four SAT valuations are
checked against their complete ground clauses.

\subsection{Constructor fusion by a two-coordinate model}
\label{se:fusion}

The absence of congruence for class constructors permits a second
reduction. Replace $H(b,t)$ everywhere by $N(b,t)$. Call the result
$\mathsf{SE19N}$, or $\mathsf{SE20N}$ for the extensional control.
The exact sizes remain 971 and 1,138 bits, but the number of proper
functions drops from nine to \textbf{eight}. We now prove both
consistency directions; identifying two functions in a presentation
does not by itself justify equiconsistency.

Every model of a fused theory expands to its unfused counterpart by
setting $H=N$. For the converse start with an inherited native
FC15$^*$ model $M$, with genuine identity as equality and the original
separate $H_M,N_M$. Put $q_M=N_M(d_M,d_M)$, so $d_M\ne q_M$.
Work on $U=M\times M$ and write $u=(u_0,u_1)$. Set
\begin{align*}
 d_U&=(d_M,d_M),&q_U&=(q_M,q_M),\\
 u\sim_U v&\ \Longleftrightarrow\ u_0=v_0,&
 E_U(u,v)&\ \Longleftrightarrow\ E_M(u_0,v_0),\\
 \mathsf{Inc}_U(s,t)&\ \Longleftrightarrow\
             \mathsf{Inc}_M(s_0,t_1),&
 G_U(b,t)&\ \Longleftrightarrow\ G_M(b_0,t_1).
\end{align*}
Define the functions by
\begin{align*}
 V_U(s,j)&=(V_M(s_0,j_0),d_M),\\
 C_U(s,j,z)&=(C_M(s_0,j_0,z_0),d_M),\\
 N_U(a,b)&=(H_M(a_0,b_1),N_M(a_1,b_1)),\\
 A_U(j,k)&=(d_M,A_M(j_0,k_0)),\\
 Q_U(j,a)&=(d_M,Q_M(j_0,a_1)),\\
 W_U(s,j,a)&=(W_M(s_0,j_0,a_1),d_M),\\
 Y_U(b,t,u)&=(Y_M(b_0,t_1,u_0),d_M),\\
 Z_U(b,t,u)&=(Z_M(b_0,t_1,u_0),d_M).
\end{align*}
Thus $N_U$ has two independent roles: its first coordinate supplies
the old hull/bad-input object and its second supplies the old NAND
class. These are compatible because incidence reads the class argument's
second coordinate while membership and the equality surrogate read
first coordinates.

Every assignment law projects to the corresponding native law.
The three NAND clauses follow from the second coordinate of $N_U$;
the atomic and quantified class clauses similarly project to $A_M,Q_M$
and $W_M$. Direct substitution gives the exact identity
\[
 R_U(b,t,x,y)=R_M(b_0,t_1,x_0,y_0).
\]
Here and below equality of displayed predicates means equality of
their truth values. The first coordinate of the derived member term is
\[
 (D_U(x))_0=
 W_M(C_M(d_M,q_M,x_0),d_M,A_M(d_M,q_M))=D_M(x_0).
\]
All six closure clauses therefore project exactly to their native
counterparts. The four selective-equality clauses hold because
$\sim_U$ is equality of first coordinates and $V_U$ respects that
coordinate in its first argument. This construction needs no
congruence for $N,A,Q,W,Y,Z$ or incidence.

For the extensional theory begin with the inherited FC16$^*$ native
model. In its additional clause, the first coordinate of every $P$
assignment is the native $P_M$, and the second coordinate of every
nested $A,N,Q$ class expression is its native class value. Its
incidence literal and three equality guards therefore reduce to the
original native extensional clause. This proves the upper direction
for $\mathsf{SE20N}$ as well. Combined with the previous lower
interpretations and inherited calibrations, ZFC proves
\[
 \operatorname{Con}(\mathsf{SE19N})\leftrightarrow
 \operatorname{Con}(\mathrm{ZFC}),\qquad
 \operatorname{Con}(\mathsf{SE20N})\leftrightarrow
 \operatorname{Con}(\mathrm{ZFC}).
\]
The independent named selector is essential to this particular
transport: $q_U$ was chosen explicitly and need not equal
$N_U(d_U,d_U)$.

There is also an optional three-predicate variant. Replace each
$G(b,t)$ by $A(b,t)\sim d$, with the same sign, and modify only the
first coordinate of $A_U$:
\[
 A_U(j,k)=
 \left(
 \begin{cases}d_M,&G_M(j_0,k_1),\\q_M,&\neg G_M(j_0,k_1),\end{cases}
 \ A_M(j_0,k_0)\right).
\]
Then $A_U(b,t)\sim_U d_U$ holds exactly when $G_M(b_0,t_1)$ does.
The atomic-class clauses still read only the unchanged second
coordinate. The lower direction defines the old $G$ by this formula.
There are six guard occurrences, so this trades the fourth predicate
for twelve term nodes, or 48 flat bits: 1,019 and 1,186 bits.

\subsection{Injective raw updates share three more witness constructors}
\label{se:witness-fusion}

Retain the primitive guard $G$ and first perform $H=N$. There is a
six-function presentation obtained by the further substitutions
\[
 Y(b,t,u):=A(C(b,t,u),d),\qquad
 Z(b,t,u):=Q(C(b,t,u),d).
\]
A five-function presentation additionally substitutes
\[
 W(s,j,a):=A(C(s,j,a),q).
\]
These substitutions apply everywhere, including inside $R$ and the
derived member term $D$, and are expanded in the emitted streams.
The respective proper signatures are $V,C,N,A,Q,W$ and $V,C,N,A,Q$.
The raw theories have 1,011 and 1,027 bits; their extensional controls
have 1,178 and 1,194. Clause and literal counts remain 19/45 and
20/49, and the maximum local variable count remains four.

Here is a full upper-model argument that rules out collisions between
the shared witness roles. In the native model $M$ the pairing
\[
 P_M(x,y)=C_M(C_M(d_M,d_M,x),q_M,y)
\]
is injective, by its two exact native projection equations. Therefore
a fixed nesting of $P_M$ gives an injection $J:M^6\to M$. Replace the
earlier definition of $C_U$ by
\[
 C_U(s,j,z)=
 \bigl(C_M(s_0,j_0,z_0),
       J(s_0,s_1,j_0,j_1,z_0,z_1)\bigr).
\]
Its first coordinate is unchanged, so all previous projection
identities and clause verifications still apply. Now $C_U:U^3\to U$
is injective. If $c$ lies in its range, let
$(\beta_c,\tau_c,\upsilon_c)$ denote its unique raw preimage.

Keep the second coordinates
\[
 (A_U(j,k))_1=A_M(j_0,k_0),\qquad
 (Q_U(j,a))_1=Q_M(j_0,a_1)
\]
for every pair of raw arguments. Define their first coordinates by
the following disjoint cases, using $d_M$ whenever no displayed case
applies:
\begin{align*}
 (A_U(c,d_U))_0
   &=Y_M((\beta_c)_0,(\tau_c)_1,(\upsilon_c)_0)
       &&(c\in\operatorname{ran}(C_U)),\\
 (Q_U(c,d_U))_0
   &=Z_M((\beta_c)_0,(\tau_c)_1,(\upsilon_c)_0)
       &&(c\in\operatorname{ran}(C_U)),\\
 (A_U(c,q_U))_0
   &=W_M((\beta_c)_0,(\tau_c)_0,(\upsilon_c)_1)
       &&(c\in\operatorname{ran}(C_U)).
\end{align*}
The first two equations handle the six-function theory, retaining
primitive $W_U$ as before. Include the third and remove $W_U$ for the
five-function theory. The two cases for $A_U$ cannot overlap:
$d_U\ne q_U$ as raw pairs. Within each case the preimage is unique.
Consequently these equations define total single-valued functions;
they impose no selector-class or raw-code congruence.

For all raw arguments the substituted witnesses have exactly the
required native first coordinates:
\begin{align*}
 (A_U(C_U(b,t,u),d_U))_0&=Y_M(b_0,t_1,u_0),\\
 (Q_U(C_U(b,t,u),d_U))_0&=Z_M(b_0,t_1,u_0),\\
 (A_U(C_U(s,j,a),q_U))_0&=W_M(s_0,j_0,a_1).
\end{align*}
Their second coordinates can take the class values already specified
by $A_U,Q_U$: every witness occurrence is consumed through its first
coordinate, either directly by $E_U$ or as a replacement argument of
$C_U$. This includes the witness inside the expanded term $D_U$.
Thus every original clause has the same truth-value projection as
before. Conversely, defining $H,Y,Z$ and optionally $W$ by the stated
terms expands any reduced model to $\mathsf{SE19}$ or
$\mathsf{SE20}$. This proves both consistency equivalences in ZFC
for each of the four new presentations.

Each replacement of $Y$ or $Z$ adds two term nodes, with five total
occurrences: ten nodes and 40 bits. Replacing $W$ adds two nodes at
each of two occurrences: four more nodes and 16 bits. These are
charged syntax costs. No injective update clause, pairing constructor,
or decoding function is added to the target. Injectivity is a property
of the model used to prove the upper direction. In the lower direction
the raw truth lemma continues to work for all models, including those
whose raw update operation is not injective.

All three retained-guard fused signatures fit a radix-eight total
ground-term codec with $d=0,q=1$ and constructor codes
$2+i+8\operatorname{tuple}(\bar n)$. The eight-function signature uses
all tags; the six- and five-function signatures have two and three
unused unary padding tags. Fair grounding uses the same polynomial
$5B^4+13B^3+1$ and its extensional $+B^3$ control. The five-function
raw target has maximum term depth four. The encoder, independent
parser, codec roundtrips and exact finite SAT witnesses are supplied
in \texttt{fused\_hull.py} and \texttt{fused\_witnesses.py}.
None of these payload sizes is a binary-machine state count.

\subsection{The six-function frozen kernel as a second compiler target}

The same method also applies to the incoming frozen-assignment kernel.
Its Boolean-normalized named-selector input \texttt{FC14Qd} has thirteen
native clauses followed by the definition $q=I(d,d)$. Delete that last
definition, interpret its other equalities by $\sim$, and append
\eqref{se:eq}--\eqref{se:sep}. The resulting frozen target has
\textbf{17 clauses, 36 literals, six functions, three predicates,
and 1,042 flat bits}. Its extensional control has 18 clauses,
40 literals and 1,201 bits. Both retain the two charged constants,
and their maximum local variable count is four. Exact inputs and
independently parsed outputs are emitted by
\texttt{frozen\_selective\_control.py}.

Here the functions are $V,C,I,A,U,F$. The raw contracts are implication
for $I$, falsity for $d$, membership atoms for $A$, and universal
quantification for $U$. The quantifier witness is the inherited term
\[
 J(s,j,a)=V(F(P(s,j),a),q).
\]
Its one witness clause still proves the exact raw universal contract:
if the universal class is false, the actual raw update at $J(s,j,a)$
is a counterexample. The truth-lemma quantifier step therefore works
just as above and assumes no congruence for this term.

The two additional native $F$ clauses give
\[
 \mathsf{Inc}(s,F(b,t))\leftrightarrow
 \mathsf{Inc}(C(C(b,d,V(s,d)),q,V(s,q)),t).
\]
For a class $t$ representing a fixed ordinary matrix with parameters
stored in $b$, the right side is exactly its relation evaluated on
$[V(s,d)],[V(s,q)]$. Thus this frozen class has precisely the two
coordinate dependencies needed by the inherited EPSC wrapper,
even when the source matrix used the wrapper's auxiliary label.
Applying the same raw truth induction to that wrapper, with an
auxiliary selector class distinct from $[d],[q]$, gives the full
supported EPSC instance on $D$. The Boolean normalization used in
the incoming input remains valid because the implication/falsity
contracts are exact on every raw assignment. Its extensional clause
gives quotient Extensionality by the same argument as $\mathsf{SE20}$.

For the upper direction, choose the existing rich-class-code native
model from the frozen-kernel note and set the new $q$ to its old
$I(d,d)$. That model has genuine equality and $d\ne q$, so all four
new completion clauses hold. Consequently the same inherited
EPSC consistency adapters prove the new frozen consistency
equivalences in ZFC. The rich-code witness-sharing theorem is
inherited; this corollary changes only equality requirements and
the selector definition.

The frozen inputs have six proper functions and three predicates.
They use fewer clauses and predicates than the five-function fused
targets, but more payload bits and substantially deeper terms: the
new frozen maximum depth is twelve. Both families admit a radix-eight
ground-term codec. Direct
grounding emits $3B^4+12B^3+B+1$ records, or $3B^4+13B^3+B+1$
for the extensional control. Neither payload family dominates all
compiler costs, and neither establishes a Turing-machine state count.

\subsection{Provenance and next compilation question}

The fifteen native clauses, guard sharing, derived member witness,
raw-EPSC interface, and extensional control come from ProveIt source44,
preserved in the reference files. The general finite assignment/class
method is attributed there to Harvey Friedman, \emph{Impossible
Counting}, Section 4.1. The preceding incoming archive supplied the
20/53 and 21/57 selective native-model reconstruction controls.
The new contribution is the raw-assignment truth lemma over an
$E$ quotient, the removal of the additional congruence obligations,
the independent-selector variant, and the exact emitted finite
payload/refuter. The two-coordinate model and injective raw-update
construction additionally justify sharing the hull, Boolean, atomic-class,
quantifier-class and witness roles in the explicit patterns proved above.
Historical novelty beyond the inspected sources is
not asserted.

The practical compiler opportunity is specific: only four local
variables remain, all update and incidence congruence checking
disappears, lookup transport compares an unchanged raw selector, and
the function signature can be reduced to five proper constructors.
A complete binary implementation must still pay for fair term
enumeration, clauses, Boolean valuation search, every helper,
and blank-tape initialization. The 971-bit stream alone gives
no numerical state bound below 245.
